\chapter{Calculator I/O \& File Formats}
\label{ch:io-types}

\newthought{Formats are a plugin concern.} A calculator's \yamlkey{input} and
\yamlkey{output} specs name a \yamlkey{type} --- \code{JSON}, \code{CSV}, \code{SLHA},
\ldots --- and the actual reading and writing is delegated to the \Portal\ package.
Upgrading \Portal\ adds formats without touching the runtime. This chapter covers the
spec anatomy, the format catalog, the \yamlkey{save} policy for keeping files, and the
\cli{Jarvis2 portal} tools.

\section{The division of labour}
\label{sec:io-ownership}

\begin{itemize}
  \item \textbf{\Jtwo\ owns the semantics}: the spec syntax, token resolution in
        \yamlkey{path}, expression evaluation in \yamlkey{Dump} variables, and the
        \yamlkey{save} policy that copies files into \file{SAMPLE/}.
  \item \textbf{\Portal\ owns the formats}: how a \code{JSON} or \code{SLHA} file is
        actually serialized and parsed --- \emph{variable read/write only}, no file
        management.
\end{itemize}

\begin{jwarn}
An unknown or empty \yamlkey{type} is a \textbf{hard failure} that lists the formats
your installed \Portal\ registers --- never a silent skip. If a format you expect is
missing from that list, upgrade \Portal.
\end{jwarn}

\section{The format catalog}
\label{sec:format-catalog}

\begin{table}[ht]
  \footnotesize
  \begin{tabular}{@{}p{0.32\linewidth}p{0.62\linewidth}@{}}
    \toprule
    \textbf{type} & \textbf{Format} \\
    \midrule
    \code{JSON}    & \textsc{json} documents; dotted \yamlkey{entry} paths navigate nesting \\
    \code{CSV} / \code{TSV} & delimited tables \\
    \code{DAT}     & whitespace-separated numeric files \\
    \code{Wolfram} & Mathematica expression files \\
    \code{SLHA} / \code{xSLHA} & \textsc{susy} Les Houches Accord (and extended) spectrum files \\
    \bottomrule
  \end{tabular}
  \caption{Formats exposed by the current \Portal\ V2 surface. Third-party packages can
  register more through the \code{jarvishep.io} entry point; the authoritative list is
  \cli{Jarvis2 portal man}.}
  \label{tab:portal-formats}
\end{table}

\section{\yamlkey{input[]}: writing calculator input}
\label{sec:input-spec}

\begin{yamlwin}
input:
  - name: lha_in
    path: "@Sdir/LesHouches.in"
    type: SLHA
    save: true                    # keep a copy under SAMPLE/
    actions:
      - type: Dump
        variables:
          - name: m0
            expression: m0            # computed (here: verbatim by expression)
          - name: tanb
            expression: "10 * tb_u"   # any shared-language expression
          - name: sign_mu             # no expression -> copy observable verbatim
\end{yamlwin}

\begin{keylist}
  \item[name] The spec's label (used in logs and generated observable names).
  \item[path] Where to write; per-Sample tokens resolve here, and \token{@Sdir}
        materializes the Sample directory.
  \item[type] The \Portal\ format.
  \item[actions] Today exactly one action type exists: \yamlkey{Dump} writes the listed
        variables into the file. Each entry writes \yamlkey{name} either from an
        \yamlkey{expression} (shared language, Chapter~\ref{ch:expressions}) or, with no
        expression, verbatim from the observable of the same name --- a missing
        observable becomes the literal string \code{MISSING\_VALUE}, which is your cue
        in the produced file.
\end{keylist}

If the target file already exists, \yamlkey{Dump} \emph{merges} into it rather than
replacing it --- useful for layering computed values onto a template your commands
copied into place first.

\section{\yamlkey{output[]}: reading calculator output}
\label{sec:output-spec}

\begin{yamlwin}
output:
  - name: spectrum
    path: "@Sdir/SPheno.spc"
    type: SLHA
    save: true
    variables:
      - name: mh
        entry: "MASS.25"          # dotted path into the parsed structure
      - name: mstop1
        entry: "MASS.1000006"
\end{yamlwin}

Each \yamlkey{variables} entry maps \yamlkey{name} from \yamlkey{entry} --- a dotted
path into the parsed file, defaulting to the name itself. A missing entry captures
\code{None} rather than failing, so downstream expressions see a null value they can be
written to tolerate (or fail on, loudly --- Section~\ref{sec:likelihood-failure}).

\section{The \yamlkey{save} policy}
\label{sec:save-policy}

By default nothing a calculator writes or reads is kept --- working files vanish with
the working directory. The \yamlkey{save} flag on a spec keeps a copy under the Sample's
directory, handled by a dedicated per-Worker file-operation process so Workers never
block on slow storage:

\begin{itemize}
  \item \textbf{input specs}: copied into \file{SAMPLE/<bucket>/<uuid>/} only when
        \yamlkey{save: true}.
  \item \textbf{output specs}: \yamlkey{save: true} copies into the Sample directory
        proper; without it, outputs that must survive the working directory go to the
        Sample's hidden \file{.temp/} area instead.
  \item Saved files keep their basename; a name that would collide across modules is
        disambiguated as \code{basename@Module}.
  \item When a spec is saved, the file's final path is also added to the observables,
        so the database records where each kept artifact lives.
\end{itemize}

\begin{jtip}
Save what you would need to re-analyse a point without re-running it --- the spectrum
file, the event-level summary --- and nothing else. At $10^6$ samples, every saved file
is a million files; buckets keep the filesystem sane
(Section~\ref{sec:sample-directory}), but nothing keeps the disk big.
\end{jtip}

\section{Portal from the command line}
\label{sec:portal-cli}

The \Portal\ registry is inspectable and drivable without a scan:

\begin{terminal}
\begin{jcmd}
Jarvis2 portal man            # list every registered format
Jarvis2 portal man slha       # one format's manual
Jarvis2 portal path/to/io.yaml   # run a standalone Portal IO YAML
\end{jcmd}
\end{terminal}

The same commands exist as the standalone \cli{jportal} tool. Running a small
\textsc{io yaml} against a real output file of your calculator is the fastest way to get
\yamlkey{entry} paths right before wiring them into a scan card.
